\input{def}

\begin{document}

%----------HEADING----------
% \begin{tabular*}{\textwidth}{l@{\extracolsep{\fill}}r}
%   \textbf{\href{http://sourabhbajaj.com/}{\Large Sourabh Bajaj}} & Email : \href{mailto:sourabh@sourabhbajaj.com}{sourabh@sourabhbajaj.com}\\
%   \href{http://sourabhbajaj.com/}{http://www.sourabhbajaj.com} & Mobile : +1-123-456-7890 \\
% \end{tabular*}

\begin{center}
  \textbf{\Huge \scshape Igor Varejão} \\ \vspace{1pt}
  \href{https://igorvarejao.tech}{\underline{igorvarejao.tech}} $|$
  \href{mailto:igor.varejao.dev@gmail.com}{\underline{igor.varejao.dev@gmail.com}} $|$ 
  \href{https://www.linkedin.com/in/igor-varejao}{\underline{Linkedin}} $|$
  \href{https://github.com/ivarejao}{\underline{Github}}
\end{center}

%-----------TEACHING EXPERIENCE-----------
\section{Experiência de Ensino}
\resumeSubHeadingListStart

\resumeSubheading
  {Monitor -- Redes Neurais e Aprendizado Profundo}{Agosto de 2025 -- Dezembro de 2025}
  {Universidade Federal de Minas Gerais}{Belo Horizonte, MG, Brasil}
  \resumeItemListStart
    \resumeItem{Atuação como monitor em um curso de Aprendizado Profundo no programa de Especialização em Ciência de Dados cobrindo Treinamento de Redes Neurais, Redes Convolucionais e Redes Neurais Recorrentes}
    \resumeItem{Condução de sessões tutoriais e horários de atendimento para ajudar estudantes a compreender conceitos complexos de aprendizado profundo e técnicas de implementação}
    \resumeItem{Fornecimento de orientação sobre implementações baseadas em PyTorch e melhores práticas para desenvolvimento de modelos de aprendizado profundo}
  \resumeItemListEnd

\resumeSubheading
  {Monitor -- Programação Competitiva (INF16169)}{Março de 2023 -- Julho de 2023}
  {Universidade Federal do Espírito Santo}{Vitória, ES, Brasil}
  \resumeItemListStart
    \resumeItem{Atuação como monitor na disciplina eletiva de Programação Competitiva no programa de Bacharelado em Ciência da Computação}
    \resumeItem{Mentoria de estudantes em técnicas de resolução de problemas algorítmicos, estruturas de dados e estratégias de otimização usando C++ e Python}
    \resumeItem{Avaliação de submissões de estudantes e fornecimento de feedback detalhado sobre eficiência de algoritmos e qualidade de código}
  \resumeItemListEnd

\resumeSubHeadingListEnd

%-----------EXPERIENCE-----------
\section{Indústria}
\resumeSubHeadingListStart

\resumeSubheading
  {Engenheiro de Dados}{Outubro de 2024 -- Maio de 2025}
  {\href{https://olhododono.agr.br}{Olho Do Dono}}{Vitória, ES, Brasil}
  \resumeItemListStart
    \resumeItem{Criação e manutenção de pipelines CI/CD, apoiando práticas DevOps para testes e implementação automatizados}
    \resumeItem{Utilização de Azure, Docker e PostgreSQL para suporte de infraestrutura e processamento de dados}
    \resumeItem{Escalonamento de sistemas de predição de peso de gado e otimização de fluxos de trabalho de visão computacional para ambientes de produção}
  \resumeItemListEnd
  
\resumeSubheading
  {Estagiário de Engenharia de Dados}{Julho de 2022 -- Setembro de 2024}
  {\href{https://olhododono.agr.br}{Olho Do Dono}}{Vitória, ES, Brasil}
  \resumeItemListStart
    \resumeItem{Desenvolvimento de sistemas de predição de peso de gado baseados em vídeos usando Python, garantindo desempenho e escalabilidade}
    \resumeItem{Realização de anotação de imagens, segmentação, marcação e ingestão de dados no CVAT com rigorosas verificações de qualidade/sanidade de dados; supervisionando uma equipe executando os mesmos fluxos de trabalho de anotação e segmentação.}
    \resumeItem{Condução de Análise Exploratória de Dados (EDA) extensiva usando Jupyter Notebook, Briefer e Plotly para avaliar e otimizar o desempenho de sistemas de IA}
    \resumeItem{Contribuição no desenvolvimento e manutenção de soluções de visão computacional e IA, coordenando múltiplos servidores através de ZMQ para comunicação eficiente.}
  \resumeItemListEnd


\resumeSubHeadingListEnd

\section{Pesquisa}
\resumeSubHeadingListStart

\resumeSubheading
  {Pesquisador de Graduação em Inteligência Artificial}{Agosto de 2020 -- Julho de 2024}
  {Universidade Federal do Espírito Santo}{Vitória, ES, Brasil}
  \resumeItemListStart
    \resumeItem{Pesquisa e desenvolvimento de redes neurais profundas para reconhecimento de padrões de defeitos em sistemas de bombeamento centrífugo submersível para o projeto DRPDBCS em parceria com a Petrobras}
    \resumeItem{Realização de análise exploratória de dados (EDA) usando Jupyter Notebook e Plotly para visualizar padrões e derivar insights acionáveis.}
    \resumeItem{Aplicação de técnicas de pré-processamento e manipulação de dados a sinais de vibração nos domínios de tempo, frequência e tempo-frequência usando Python (NumPy, Pandas)}
    \resumeItem{Desenvolvimento de um framework público para centralizar conjuntos de dados públicos usando Python para pesquisa em aprendizado profundo no reconhecimento de padrões de defeitos no \href{https://github.com/ivarejao/vibdata.git}{repositório github vibdata}.}
    \resumeItem{Treinamento de redes neurais profundas usando os frameworks pytorch, pytorch-lightning e scikit-learning}
    \resumeItem{Condução de testes estatísticos abrangentes (por exemplo, teste-t, teste-z, teste de Wilcoxon, Welch) para avaliar rigorosamente a significância e validade dos resultados da pesquisa.}
    \resumeItem{Autoria e publicação de artigo intitulado \href{https://doi.org/10.1016/j.ymssp.2025.112822}{\textit{The similarity bias problem: What it is and how it impacts vibration based intelligent fault diagnosis}} como resultado desta pesquisa.}

  \resumeItemListEnd


\resumeSubheading
  {Pesquisador Voluntário em Inteligência Artificial}{Agosto de 2024 -- Dezembro de 2024}
  {Universidade Federal do Espírito Santo}{Vitória, ES, Brasil}
  \resumeItemListStart
    \resumeItem{Desenvolvimento de modelos de redes neurais profundas usando aprendizado por transferência para diagnóstico de falhas em rolamentos de máquinas rotativas, em colaboração com o Departamento de Mecânica da Universidade de Ottawa, Canadá.}
    \resumeItem{Construção de um framework de benchmarking para avaliar um modelo pré-treinado em conjuntos de dados públicos centralizados de rolamentos, comparando seu desempenho com abordagens de pré-treinamento de última geração.}
    \resumeItem{Publicação de artigo de pesquisa \href{https://doi.org/10.36001/ijphm.2025.v16i3.4239 }{\textit{Towards a Universal Vibration Analysis Dataset A Framework for Transfer Learning in Predictive Maintenance and Structural Health Monitoring}}}
  \resumeItemListEnd

  
  
\resumeSubHeadingListEnd

%-----------EDUCATION-----------
\section{Educação}
  \resumeSubHeadingListStart
    \resumeSubheading
      {Universidade Federal de Minas Gerais}{Belo Horizonte, MG, Brasil}
      {Mestrado em Ciências da Computação, Inteligência Artificial}{Agosto de 2025 -- Presente}
      \resumeItemListStart
        \resumeItem{Foco de pesquisa atual: Aprendizado de Máquina para Geração Musical}
      \resumeItemListEnd
    \resumeSubheading
      {Universidade Federal do Espírito Santo}{Vitória, ES, Brasil}
      {Bacharelado em Ciência da Computação}{Março de 2019 -- Novembro de 2024}
  \resumeSubHeadingListEnd


%-----------PROJECTS-----------
\section{Projetos Pessoais}
\resumeSubHeadingListStart
\resumeProjectHeading
    {\textbf{News WebScraper} \href{https://github.com/ivarejao/veja-news-scraper}{(Github)} $|$ \emph{Python, Selenium} }{Maio de 2022 -- Julho de 2023}
    \resumeItemListStart
        \resumeItem{Desenvolvimento de uma ferramenta de web scraping para coleta de corpus textual de um site de jornal para o estudo de neologismos na cultura digital}
        \resumeItem{Implementação do Selenium para automatizar navegação de páginas e extração de conteúdo web}
        \resumeItem{Utilização do BeautifulSoup para analisar, limpar e estruturar os dados textuais para construção de corpus}
    \resumeItemListEnd
\resumeSubHeadingListEnd



%-------------------------------------------
\end{document}
